\chapter{电子与流行：用层次做质感，不用堆音色}
\label{ch:electronic}

\begin{quote}
\itshape 电子编曲的“高级感”很少来自多买几个音色，而来自每一层只做一件事：低频负责重量，主体负责身份，空气负责细节，节奏负责推进。
\end{quote}

\section{电子音色不是一件乐器，是一叠层}
\label{sec:electronic-layers}

流行与电子制作里，一个“合成器音色”往往由几层组成。把它们拆开写，比在一个音色里塞满参数更容易控制：

\begin{itemize}
\item \textbf{低频层}：负责 Sub 与重量，通常在最低频段；它不需要旋律复杂，只需要稳定、准、不抖。
\item \textbf{主体层}：负责音色的“身份”，比如副歌的宽 Pad、主奏的 Pluck 或 Lead；它决定听众记住的是什么颜色。
\item \textbf{空气层}：负责高频细节与光泽，比如短促的 High Pluck、噪声质感、句尾的亮片；它不能持续轰鸣。
\item \textbf{节奏层}：负责瞬态与推进，比如电子鼓、短促的 Stab、侧链感的留空；它决定段落有没有“走起来”。
\end{itemize}

这四层不是四条固定轨道，而是四种功能。一个音色可以同时承担两层，但不要让同一层由五个音色重复；那会变成“很响”，不是“很厚”。

\section{合成器参数：只调对声音有影响的那几个}
\label{sec:synth-params}

合成器面板很长，编曲阶段先抓住四个参数：

\begin{enumerate}
\item \textbf{包络}：Attack 决定音头是“啪”还是“渐入”，Decay/Release 决定尾巴多长；短 Attack 适合 Pluck/Stab，长 Attack 适合 Pad。
\item \textbf{滤波器}：Cutoff 决定亮暗，Resonance 决定峰有多尖；主体层不要和低频层抢同一个 cutoff，否则两层会糊成一条。
\item \textbf{音高/振荡器}：Sub 与 Bass 用低八度，Pad 用中音区，Pluck 用高八度；不要让所有层都从同一个八度出发。
\item \textbf{调制}：LFO 或包络调制可以让 Pad 有呼吸、让 Pluck 有微微的移动；调制幅度过大，会变成“抖”，不是“活”。
\end{enumerate}

参数不是越复杂越好。一个 Pad 只要 Attack 慢、Cutoff 中低、Release 长，就已经能铺住副歌；把五个旋钮都拧满，通常只会得到一坨没有焦点的声音。

\section{频段分工：低频、主体、空气各占一块}
\label{sec:frequency-roles}

\term{频段分工}{frequency-band-division}是电子编曲不糊的关键：

\begin{longtable}{p{0.18\textwidth}p{0.26\textwidth}p{0.42\textwidth}}
\toprule
层 & 频段/音区 & 写作要点 \\
\midrule
\endfirsthead
\toprule
层 & 频段/音区 & 写作要点 \\
\midrule
\endhead
低频层 & 最低频段 & Sub 与 Bass 只留一个主导；另一个让路或只落强拍 \\
主体层 & 中频 & Pad/Lead/Pluck 选一个当身份；不要三件同时占满中频 \\
空气层 & 高频 & 短句、句尾、装饰；持续高频会刺耳，也会盖住人声齿音 \\
节奏层 & 全频段的瞬态 & 鼓与 Stab 负责“点”，不要把每个十六分都填满 \\
\bottomrule
\end{longtable}

低频层最容易出问题：Sub、Bass、Kick 三件都想占最低处，结果谁都不清楚。解决办法不是加压缩，而是编排上分工——Kick 管瞬态，Bass 管根音，Sub 只在长音段落补重量。

\section{音色分层：让每层有独立的表情}
\label{sec:timbre-layering}

\term{音色分层}{timbre-layering}不是把两个 Preset 叠在一起，而是让每层有独立的时间形状：

\begin{itemize}
\item \textbf{短 Attack＋短 Release}：Pluck、Stab、电子鼓，负责“点”。
\item \textbf{长 Attack＋长 Release}：Pad、铺底，负责“面”。
\item \textbf{中 Attack＋长 Release}：Lead、Bass，负责“线”。
\item \textbf{噪声/高频短句}：Air，负责“亮”。
\end{itemize}

同一段里，“点、面、线、亮”各有其一，耳朵才不会累。如果副歌里所有层都是长音，听起来会像一团棉花；如果所有层都是短点，听起来会像机关枪。

\section{采样与循环：素材是材料，不是编曲}
\label{sec:sample-loop}

采样与循环能快速建立质感，但它们不能替你做结构：

\begin{itemize}
\item \textbf{采样}：把导入的音频放进\term{素材池}{asset pool}，再摆到音频轨；先确认长度、采样率与节拍对得上，不要靠拉伸硬救。
\item \textbf{循环}：用复制片段（Ctrl+C / Ctrl+V）拼出重复段落，比在一个长片段里硬循环更容易控制进出；每次复制后改一处，不要整段照搬。
\item \textbf{素材让路}：采样循环如果已经含鼓/贝斯，就不要再叠一套同功能的 MIDI；否则低频和瞬态会打架。
\item \textbf{版权与来源}：只用你有权使用的素材；回执与工程里保留来源，不把“网上下的”当默认可用。
\end{itemize}

循环段落最怕“从头到尾一个 loop”。解决方法是把 loop 当节奏骨架，在上面加一两个 MIDI 层做变化：副歌换 Pluck，桥段抽掉鼓，尾奏留 Pad。

\section{流行歌的电子层：主歌省，副歌给}
\label{sec:pop-layers}

流行歌的电子编曲，最常见的能量曲线是：

\begin{enumerate}
\item \textbf{主歌}：鼓＋Bass＋一件主体音色（Pad 或 Pluck）＋人声；空气层只在句尾点一下。
\item \textbf{预副歌}：加一层上升的 Pad 或 Filter 打开，制造“要来了”的感觉；不要提前把副歌所有层开满。
\item \textbf{副歌}：低频、主体、空气、节奏四层全开，但每层仍有分工；Lead 与人声错开音区。
\item \textbf{桥段}：抽掉鼓或低频，留 Pad/采样，让耳朵休息；再把副歌的层逐层加回来。
\end{enumerate}

\textbf{副歌不是“把所有轨都打开”，而是“每层都更清楚”。} 如果副歌只是音量更大、轨更多，听众会觉得吵；如果每层各就各位，即使轨数一样，也会觉得“满了”。

\begin{figure}[htbp]
  \centering
  \begin{tikzpicture}[font=\small]
    \draw[->] (0,0) -- (9.6,0) node[right] {bars};
    \foreach \x in {1,...,8} {\draw[gray!30] (\x,0.1) -- (\x,3.8); \node[font=\scriptsize] at (\x,0.05) {\x};}
    \fill[yebanblue!12] (0.2,0.2) rectangle (4.2,3.6);
    \fill[yebanblue!18] (4.2,0.2) rectangle (8.2,3.6);
    \node[font=\scriptsize] at (2.2,3.4) {verse};
    \node[font=\scriptsize] at (6.2,3.4) {chorus};
    \fill[gray!45] (0.2,2.95) rectangle (8.2,3.25);
    \fill[yebanblue!55] (0.2,2.15) rectangle (4.2,2.55);
    \fill[yebanblue!75] (4.2,2.15) rectangle (8.2,2.55);
    \fill[orange!45] (0.2,1.35) rectangle (4.2,1.75);
    \fill[orange!70] (4.2,1.35) rectangle (8.2,1.75);
    \fill[green!35] (0.2,0.55) rectangle (4.2,0.95);
    \fill[green!55] (4.2,0.55) rectangle (8.2,0.95);
    \node[font=\scriptsize, anchor=west] at (8.4,3.1) {drums};
    \node[font=\scriptsize, anchor=west] at (8.4,2.35) {body/pad};
    \node[font=\scriptsize, anchor=west] at (8.4,1.55) {bass/sub};
    \node[font=\scriptsize, anchor=west] at (8.4,0.75) {air/pluck};
  \end{tikzpicture}
  \caption{主歌到副歌的四层能量：副歌不是多一条轨，而是每层都更完整；空气层只在副歌持续。}
  \label{fig:electronic-layers}
\end{figure}

\section{在夜半里落地}
\label{sec:electronic-yeban}

\begin{enumerate}
\item \textbf{建 MIDI 合成器轨}：用「轨道 → + 新建 MIDI 合成器轨（Ctrl+N）」建轨，按低频、主体、空气、节奏命名；一条轨只承担一层。
\item \textbf{在设备机架里调音色}：在底部设备机架或左侧检查器里调合成器参数；先把 Attack/Release 与 Cutoff 调对，再谈其他旋钮。
\item \textbf{用自动化泳道做段落变化}：在编曲区展开自动化泳道（A），把 Filter Cutoff、音量等参数随时间写进工程；副歌打开、主歌收住，不靠手推。
\item \textbf{导入采样与循环}：用「文件 → 导入音频文件...」把素材导入素材池，再摆到音频轨；采样轨与 MIDI 轨分开，方便 Solo/Mute 试听。
\item \textbf{用复制片段拼循环}：选中片段后 Ctrl+C / Ctrl+V，再拖动到下一段；复制后改一处，不要整段一模一样。
\item \textbf{用力度泳道做瞬态}：Pluck/Stab 的第一个音给足力度，Pad 保持平稳；不要让所有层都用同一个力度曲线。
\item \textbf{用 Solo/Mute 验收}：先 Solo 低频层，确认 Sub/Bass/Kick 不打架；再 Solo 主体层，确认它没有盖住人声；最后全开，确认空气层只在句尾出现。
\end{enumerate}

\screenshot{arrangement-electronic-layers}{低频、主体、空气、节奏四条 MIDI 轨与一条采样轨同屏，钢琴卷帘与自动化泳道展示副歌进出的层次}

\section{常见错误与修正}
\label{sec:electronic-mistakes}

\begin{longtable}{p{0.34\textwidth}p{0.58\textwidth}}
\toprule
常见错误 & 修正 \\
\midrule
\endfirsthead
\toprule
常见错误 & 修正 \\
\midrule
\endhead
所有合成器都占满中频 & 按低频/主体/空气分工；主体层只留一个身份。 \\
Sub、Bass、Kick 三层抢最低处 & Kick 管瞬态，Bass 管根音，Sub 只补长音重量。 \\
所有层都是长 Attack/长 Release & 加入短 Attack 的 Pluck/Stab，做“点”；面和点要并存。 \\
副歌只是把轨全部打开 & 每层更完整、更清楚，而不是更多轨；空气层只在副歌持续。 \\
采样 loop 替你做结构 & loop 当骨架，在上面加 MIDI 变化；不要整首一个 loop。 \\
所有层的 Cutoff/包络一样 & 主体层与低频层用不同 Cutoff；Pad 与 Pluck 用不同包络。 \\
Pluck/Stab 太密盖住人声 & 只留关键拍，句尾/空拍再进；不要每个十六分都点。 \\
\bottomrule
\end{longtable}

\section{练习}
\label{sec:electronic-exercise}

\begin{exercise}
用八小节做一段电子流行副歌：
\begin{enumerate}
\item 先建四条轨：低频、主体 Pad、空气 Pluck、电子鼓；每轨只做一层。
\item 低频层写根音，电子鼓写四拍骨架；确认 Kick 与 Bass 不在同一拍点抢最低处。
\item 主体 Pad 用长 Attack/长 Release，只占中频；空气 Pluck 只写句尾两拍，不要连续。
\item 用自动化泳道让 Pad 的 Cutoff 在第 5 小节打开；把 Pluck 的力度做成“第一拍强、后面弱”。
\end{enumerate}
最后 Solo 人声：副歌的电子层有没有盖住歌词？如果拿掉空气层更清楚，说明它已经完成了任务。
\end{exercise}

\section{本章要点}
\label{sec:electronic-takeaways}

\begin{itemize}
\item 电子音色是低频、主体、空气、节奏四层功能，不是四个 Preset。
\item 编曲阶段只调包络、滤波器、音高与调制；参数服务音乐，不是参数比赛。
\item 低频、主体、空气各占频段；Sub/Bass/Kick 先分工，再谈压缩。
\item 音色分层要有“点、面、线、亮”，不要全是长音或全是短点。
\item 采样与循环是材料；结构仍要自己写，loop 之上要加变化。
\item 在夜半里用合成器轨、设备机架、自动化、素材池、复制片段与 Solo/Mute 落地。
\end{itemize}

\begin{helpmapping}
\begin{longtable}{p{0.45\textwidth}p{0.45\textwidth}}
\toprule
小节 & 帮助锚点 \\
\midrule
电子音色不是一件乐器 & \texttt{arrangement-electronic-layers} \\
合成器参数 & \texttt{arrangement-synth-params} \\
频段分工 & \texttt{arrangement-frequency-roles} \\
音色分层 & \texttt{arrangement-timbre-layering} \\
采样与循环 & \texttt{arrangement-sample-loop} \\
流行歌的电子层 & \texttt{arrangement-pop-layers} \\
在夜半里落地 & \texttt{arrangement-electronic-yeban} \\
常见错误与修正 & \texttt{arrangement-electronic-mistakes} \\
练习 & \texttt{arrangement-electronic-exercise} \\
\bottomrule
\end{longtable}
\end{helpmapping}
